\chapter{KHẢO SÁT HIỆN TRẠNG VÀ YÊU CẦU HỆ THỐNG}

\section{Khảo sát}

Theo Báo cáo Thị trường Bất động sản TP. Hồ Chí Minh 2025, lượng người
lao động ngoại tỉnh, sinh viên và chuyên gia trẻ thuê trọ tại các đô thị
lớn vượt 4,5 triệu người. Mô hình kinh doanh nhà trọ phân khúc tầm trung
(3-15 triệu đồng/phòng/tháng) đã trở thành nguồn thu nhập chính của hàng
chục nghìn hộ kinh doanh cá thể, trong đó nhiều chủ trọ sở hữu chuỗi
3-15 cơ sở phân tán ở nhiều quận/huyện khác nhau.

Trong quá trình khảo sát thực địa và phỏng vấn sâu 25 chủ trọ và hơn 200
khách thuê tại địa bàn TP. HCM, Bình Dương và Đồng Nai, nhóm thực hiện
đồ án đã nhận diện được các vấn đề trọng yếu sau:

\begin{itemize}
\item
  Việc quản lý đa cơ sở (chuỗi nhà trọ) gần như không thực hiện được nếu
  chỉ dùng giấy tờ hoặc Excel rời rạc - chủ trọ không có cái nhìn tổng
  quan về doanh thu, tỉ lệ lấp đầy và công nợ giữa các cơ sở.
\item
  Sai sót trong khâu chốt chỉ số điện nước hàng tháng (gần 30\% chủ trọ
  thừa nhận từng tính sai ≥ 1 lần/tháng), dẫn đến tranh chấp với khách
  thuê và mất uy tín.
\item
  Quy trình lập hợp đồng, gia hạn, chấm dứt hợp đồng đều thủ công, mất
  nhiều thời gian và dễ thất lạc giấy tờ.
\item
  Khách thuê không có kênh truy cập thông tin hoá đơn, lịch sử thanh
  toán một cách minh bạch và thuận tiện; phần lớn nhận hoá đơn dưới dạng
  tin nhắn Telegram hoặc ảnh chụp viết tay.
\item
  Việc thu - chi tiền mặt là chủ yếu, dẫn đến công nợ phát sinh và khó
  kiểm soát; ít chủ trọ áp dụng được cổng thanh toán điện tử quy mô lớn.
\item
  Khách vãng lai tìm phòng phải tự gọi điện từng khu trọ, không có kênh
  trực tuyến để xem chi tiết phòng còn trống, đặt cọc giữ phòng nhanh
  chóng.
\item
  Việc đánh giá hiệu quả vận hành của từng quản lý không có dữ liệu định
  lượng để đối chiếu.
\end{itemize}

Các vấn đề trên cho thấy tính cấp thiết của một hệ thống thông tin quản
lý chuỗi nhà trọ ở quy mô doanh nghiệp vừa, giúp các chủ trọ chuyển từ
mô hình "quản lý hộ kinh doanh" sang mô hình "vận hành dịch vụ chuyên
nghiệp".

\section{Yêu cầu về chức năng hệ thống}

\subsection{Yêu cầu về chức năng}

Hệ thống cần đáp ứng đầy đủ các nhóm chức năng nghiệp vụ được tổng hợp
từ quá trình khảo sát:

\begin{longtable}[]{@{}>{\RaggedRight\arraybackslash}p{0.27\linewidth}>{\RaggedRight\arraybackslash}p{0.68\linewidth}@{}}
\toprule
\textbf{Nhóm chức năng} & \textbf{Mô tả chi tiết}\tabularnewline
\midrule
\endhead
Quản lý xác thực \& tài khoản & Đăng ký, đăng nhập (email + mật khẩu),
đăng xuất, quên/đặt lại mật khẩu, cập nhật hồ sơ, khoá/mở khoá tài
khoản, phân quyền theo vai trò (Chủ trọ kiêm Admin, Quản lý, Khách thuê,
Khách vãng lai).\tabularnewline
Quản lý chuỗi nhà trọ - chi nhánh & Thêm/sửa/ngừng hoạt động nhà trọ;
gán quản lý; thiết lập địa chỉ, toạ độ, ảnh đại diện, thông tin liên hệ;
thống kê đa cơ sở.\tabularnewline
Quản lý phòng \& tài sản & CRUD loại phòng (giá cơ bản, diện tích, tiện
nghi), CRUD phòng (mã, tầng, giá, trạng thái: trống/đang thuê/đặt cọc),
quản lý tài sản trong phòng và tình trạng tài sản.\tabularnewline
Quản lý hợp đồng \& khách thuê & Tạo hồ sơ khách thuê, ký hợp đồng (có
ký số trực tuyến), gia hạn, sửa đổi, chấm dứt hợp đồng.\tabularnewline
Dịch vụ \& đơn giá & Cấu hình các dịch vụ tính theo chỉ số (điện, nước),
tính theo khoán (internet, vệ sinh, gửi xe, an ninh) và áp dụng đơn giá
bậc thang.\tabularnewline
Ghi chỉ số \& tính tiền & Ghi nhận chỉ số điện nước hàng tháng, tự động
tính tiền tiêu thụ và tổng hợp các khoản phí cố định.\tabularnewline
Quản lý hoá đơn & Tạo hoá đơn (thủ công/tự động), gửi hoá đơn qua email,
nhắc thanh toán tự động, theo dõi trạng thái và lịch sử thanh
toán.\tabularnewline
Thanh toán & Hỗ trợ thanh toán online qua cổng VNPay Sandbox (tích hợp
thật) và mã VietQR động; xác nhận đã thu tiền mặt (luồng pending\_cash);
phát hành biên lai điện tử.\tabularnewline
Báo cáo \& thống kê & Dashboard tổng quan; báo cáo doanh thu theo
tháng/quý/năm; tỉ lệ lấp đầy; công nợ khách thuê; chi phí vận hành; xuất
Excel/PDF.\tabularnewline
Thông báo \& nhắc nhở & Đẩy thông báo qua email và in-app (trong ứng
dụng) về: hợp đồng sắp hết hạn, hoá đơn đến hạn, trạng thái thay
đổi.\tabularnewline
Tìm kiếm \& đặt cọc phòng & Khách vãng lai tìm kiếm phòng theo khu vực,
giá, tiện nghi; đặt cọc giữ phòng online.\tabularnewline
\bottomrule
\end{longtable}

\subsection{Yêu cầu về phi chức năng}

\begin{longtable}[]{@{}>{\RaggedRight\arraybackslash}p{0.27\linewidth}>{\RaggedRight\arraybackslash}p{0.68\linewidth}@{}}
\toprule
\textbf{Tiêu chí} & \textbf{Yêu cầu}\tabularnewline
\midrule
\endhead
Hiệu năng & Thời gian phản hồi trung bình dưới 1.5 giây cho các thao tác
truy vấn thông thường; hỗ trợ tối thiểu 500 người dùng đồng thời; xử lý
lô tạo hoá đơn cho 5.000 phòng trong vòng 5 phút.\tabularnewline
Bảo mật & Mật khẩu lưu dưới dạng bcrypt; xác thực OTP qua email (khi
đăng ký \& khôi phục mật khẩu); phân quyền theo vai trò (RBAC); ghi nhật
ký toàn bộ hành động nhạy cảm; tuân thủ Nghị định 13/2023/NĐ-CP về bảo
vệ dữ liệu cá nhân.\tabularnewline
Khả dụng & SLA 99.5\%; cơ chế sao lưu dữ liệu hàng ngày; khôi phục thảm
hoạ dưới 4 giờ (RTO ≤ 4h, RPO ≤ 1h).\tabularnewline
Tiện dụng & Giao diện đáp ứng (responsive) cho desktop, tablet và
mobile; thao tác chính hoàn tất trong ≤ 3 bước nhấp chuột; hỗ trợ song
ngữ Việt - Anh.\tabularnewline
Khả năng mở rộng & Kiến trúc đa người dùng (multi-tenant); cho phép tích
hợp module mới (chấm công nhân viên, IoT công tơ điện, AI dự báo lấp
đầy) thông qua REST API.\tabularnewline
Tương thích & Tương thích các trình duyệt hiện đại (Chrome 90+, Edge
90+, Safari 14+); hỗ trợ Android 8+ và iOS 13+ cho ứng dụng di động đi
kèm.\tabularnewline
Kiểm thử \& chất lượng & Bao phủ kiểm thử tự động (unit + integration)
tối thiểu 70\%; quy trình CI/CD tự động trên mỗi pull
request.\tabularnewline
Tuân thủ pháp lý & Lưu hoá đơn điện tử theo Nghị định
123/2020/NĐ-CP.\tabularnewline
\bottomrule
\end{longtable}

\section{Các biểu mẫu phỏng vấn}

Quá trình thu thập yêu cầu sử dụng kết hợp ba phương pháp: (i) phỏng vấn
sâu nhóm chủ trọ và quản lý chi nhánh; (ii) khảo sát online dành cho
khách thuê; (iii) quan sát thực địa quy trình nghiệp vụ tại 5 cơ sở thí
điểm. Dưới đây là tổng hợp các nhóm câu hỏi chính:

\begin{longtable}[]{@{}>{\RaggedRight\arraybackslash}p{0.20\linewidth}>{\RaggedRight\arraybackslash}p{0.37\linewidth}>{\RaggedRight\arraybackslash}p{0.37\linewidth}@{}}
\toprule
\textbf{Đối tượng} & \textbf{Câu hỏi phỏng vấn} & \textbf{Mục
đích}\tabularnewline
\midrule
\endhead
\begin{minipage}[t]{\linewidth}\raggedright
Chủ trọ (Admin)\strut
\end{minipage} & \begin{minipage}[t]{\linewidth}\raggedright
1. Hiện tại anh/chị đang quản lý bao nhiêu cơ sở và bao nhiêu phòng?
Bằng công cụ gì?

2. Thách thức lớn nhất khi tổng hợp doanh thu, công nợ và đối soát hoá
đơn giữa các cơ sở là gì?

3. Anh/chị có giao một phần công việc cho quản lý vận hành không? Việc
giám sát họ ra sao?

4. Tỉ lệ khách thuê thanh toán bằng tiền mặt và chuyển khoản là bao
nhiêu? Anh/chị có muốn áp dụng cổng thanh toán online không?

5. Anh/chị mong muốn báo cáo dạng nào (Excel, PDF, Dashboard)? Tần suất
ra sao?\strut
\end{minipage} & \begin{minipage}[t]{\linewidth}\raggedright
Xác định pain-point ở quy mô doanh nghiệp; thu thập yêu cầu phân quyền,
hoá đơn, công nợ, báo cáo.\strut
\end{minipage}\tabularnewline
\begin{minipage}[t]{\linewidth}\raggedright
Quản lý\strut
\end{minipage} & \begin{minipage}[t]{\linewidth}\raggedright
1. Anh/chị mất bao nhiêu thời gian/ngày cho việc ghi chỉ số điện nước \&
tạo hoá đơn?

2. Quy trình xác nhận thu tiền mặt và đối soát với chủ trọ hiện tại như
thế nào?

3. Anh/chị có cần ứng dụng di động để chốt số ngay tại phòng và xác nhận
thanh toán không?\strut
\end{minipage} & \begin{minipage}[t]{\linewidth}\raggedright
Thu thập yêu cầu thao tác di động, xác nhận thu tiền mặt và quy trình
vận hành hằng ngày.\strut
\end{minipage}\tabularnewline
\begin{minipage}[t]{\linewidth}\raggedright
Khách thuê\strut
\end{minipage} & \begin{minipage}[t]{\linewidth}\raggedright
1. Bạn muốn nhận thông báo hoá đơn qua kênh nào (email, in-app, push)?

2. Bạn có cần xem lại lịch sử thanh toán 6-12 tháng gần nhất không?

3. Bạn sẵn sàng thanh toán online qua VNPay/MoMo/QR Banking không?

4. Bạn có muốn xem hợp đồng đã ký dạng số (PDF) và lưu trữ trực tuyến
không?\strut
\end{minipage} & \begin{minipage}[t]{\linewidth}\raggedright
Xây dựng trải nghiệm khách thuê: tra cứu hợp đồng - hoá đơn và thanh
toán.\strut
\end{minipage}\tabularnewline
\begin{minipage}[t]{\linewidth}\raggedright
Khách vãng lai\strut
\end{minipage} & \begin{minipage}[t]{\linewidth}\raggedright
1. Bạn thường tìm phòng trọ qua kênh nào (Facebook, Chợ Tốt,
website\ldots)?

2. Bạn quan tâm thông tin gì nhất khi xem một phòng?

3. Bạn sẵn sàng đặt cọc giữ phòng online không?\strut
\end{minipage} & \begin{minipage}[t]{\linewidth}\raggedright
Phát triển chức năng tìm kiếm và đặt cọc phòng.\strut
\end{minipage}\tabularnewline
\bottomrule
\end{longtable}

